% Generated draft: requires mathematical and source review.
\documentclass[11pt]{article}
\usepackage[margin=25mm]{geometry}
\usepackage{fontspec}
\setmainfont{Latin Modern Roman}
\usepackage{amsmath,amssymb,mathtools,unicode-math}
\setmathfont{Latin Modern Math}
\usepackage{enumitem,needspace,xurl,hyperref}
\setlength{\parindent}{0pt}
\setlength{\parskip}{5pt}
\setlength{\emergencystretch}{3em}
\newcommand{\N}{\mathbb{N}}
\newcommand{\Z}{\mathbb{Z}}
\newcommand{\Q}{\mathbb{Q}}
\newcommand{\R}{\mathbb{R}}
\newcommand{\C}{\mathbb{C}}
\newcommand{\F}{\mathbb{F}}
\newcommand{\A}{\mathbb{A}}
\newcommand{\PP}{\mathbb{P}}
\newcommand{\E}{\mathbb{E}}
\newcommand{\eps}{\varepsilon}
\newcommand{\dd}{\,\mathrm{d}}
\begin{document}
\section*{UnsolvedMath ID: 2100503}
\textbf{Source problem code:} catalogue entry
\subsection{Definitions and mathematical statement}
A Poisson structure is a bivector defining a skew bracket $\{f,g\}$ on functions that satisfies the Leibniz rule and Jacobi identity. Two structures $P_0,P_1$ are compatible if every linear combination
\[
 P_{\lambda,\mu}=\lambda P_0+\mu P_1
\]
is Poisson. At a point, these give a pencil of skew-symmetric bilinear forms on the cotangent space. Its Jordan--Kronecker data consist of its Jordan-block and Kronecker-block invariants, including eigenvalue and block-size information, with details depending on the coefficient field and equivalence convention.

The supplied assessment interprets ``algebraic type'' as these pencil invariants and asks whether every such admissible type is realized by a bi-integrable dynamical system, or a complete commuting Hamiltonian family for the pencil. A common-involution requirement for a candidate family $f_1,\ldots,f_r$ would be
\[
 \{f_i,f_j\}_{P_0}=\{f_i,f_j\}_{P_1}=0\qquad\text{for all }i,j,
\]
which implies involution for all pencil members. This condition alone does not specify completeness, independence, common Hamiltonian dynamics, or nontriviality, so it is not substituted for the source's undefined bi-integrability requirement.

The literal target is existence for \emph{each algebraic type}. To turn it into a single realization theorem one must specify the admissible Jordan--Kronecker types, real/complex field, changes of pencil parameters allowed in the type equivalence, linear versus manifold setting, local versus global existence, and whether type is prescribed at one point or generically. One must also fix precisely what a bi-integrable system means, including completeness and treatment of singular leaves. Those choices are absent from the source record, so no unrestricted matrix-pencil realization claim or vacuous zero-dynamics example is silently adopted.
\subsection{Short English statement}
Does every allowable Jordan--Kronecker type of a compatible Poisson pencil support a genuinely bi-integrable system? The supplied assessment describes the intended pencil interpretation, cites broad example-producing theories, and labels the full existence problem open/unverified. It explicitly says no resolving citation was verified in that session. This is a limitation of the supplied review, not proof that the question is currently unresolved in every possible formulation.
\subsection{Sources}
\begin{itemize}
\item UnsolvedMath permanent ID 2100503, source code catalogue entry, ``Open Problems in Integrable Systems --- Bi-Poisson vector spaces.'' Registry curation and metadata: Ulam AI and UnsolvedMath Contributors, CC BY 4.0; updated 2026-08-06.
\item \emph{Open Problems, Questions, and Challenges in Finite-Dimensional Integrable Systems} (2018), Problem 5.3: \url{https://arxiv.org/abs/1804.03737}. Supplied source-TeX extraction accessed 2026-07-29. The metadata says the source presents an open item and marks status/release and rights review NEEDS\_REVIEW.
\item The supplied summary names Magri, Gel'fand--Zakharevich, and Bolsinov--Zhang for bi-Hamiltonian theory and examples, without complete titles, dates, theorem statements or locators. No missing bibliography is inferred.
\end{itemize}
\end{document}
